\documentclass[12pt,a4paper]{article}

% -------------------------------------------------
% PAQUETES
% -------------------------------------------------
\usepackage[spanish]{babel}
\usepackage[utf8]{inputenc}
\usepackage[T1]{fontenc}
\usepackage[a4paper,margin=2.5cm]{geometry}
\usepackage{setspace}
\usepackage{parskip}
\usepackage{amsmath,amssymb}
\usepackage{graphicx}
\usepackage{float}
\usepackage{array}
\usepackage{longtable}
\usepackage{booktabs}
\usepackage{xcolor}
\usepackage{hyperref}
\usepackage{enumitem}
\usepackage{listings}
\usepackage{titlesec}
\usepackage{fancyhdr}

% -------------------------------------------------
% CONFIGURACIÓN
% -------------------------------------------------
\hypersetup{
    colorlinks=true,
    linkcolor=blue,
    urlcolor=blue,
    citecolor=blue
}

\setstretch{1.15}

\pagestyle{fancy}
\fancyhf{}
\fancyhead[L]{AOC2 - Proyecto 1}
\fancyhead[R]{Memoria}
\fancyfoot[C]{\thepage}

\titleformat{\section}{\large\bfseries}{\thesection.}{0.5em}{}
\titleformat{\subsection}{\normalsize\bfseries}{\thesubsection.}{0.5em}{}

\lstset{
    basicstyle=\ttfamily\small,
    breaklines=true,
    frame=single,
    columns=fullflexible,
    keepspaces=true
}

% -------------------------------------------------
% COMANDOS ÚTILES
% -------------------------------------------------
\newcommand{\campo}[1]{\textbf{#1}}
\newcommand{\pendiente}{\textcolor{red}{[Rellenar]}}
\newcommand{\codigo}[1]{\texttt{#1}}

% -------------------------------------------------
% DOCUMENTO
% -------------------------------------------------
\begin{document}

% -------------------------------------------------
% PORTADA
% -------------------------------------------------
\begin{titlepage}
    \centering

    {\Large Universidad de Zaragoza \par}
    \vspace{0.3cm}
    {\Large EINA \par}
    \vspace{1.2cm}

    {\Huge \textbf{AOC2} \par}
    \vspace{0.5cm}
    {\LARGE Proyecto 1 \par}
    \vspace{0.3cm}
    {\Large Gestión de riesgos, excepciones y ALU multiciclo en MIPS \par}

    \vfill

    \begin{flushleft}
        \campo{Apellidos y nombre:} Tahir Berga Celma-Pablo Villa Camañes \\[0.3cm]
        \campo{Fecha:} \pendiente
    \end{flushleft}

    \vfill
\end{titlepage}

\tableofcontents
\newpage

% -------------------------------------------------
\section{Breve resumen}
% -------------------------------------------------

Redactar aquí un resumen general del trabajo realizado.
La plantilla nueva indica una extensión orientativa de unas 200 palabras.

\vspace{0.5cm}
\noindent\pendiente

% -------------------------------------------------
\section{Verificación de los requisitos funcionales}
% -------------------------------------------------

En cada requisito funcional conviene incluir:
\begin{itemize}[leftmargin=1.2cm]
    \item breve título indicando el objetivo del programa de prueba o testbench,
    \item código ensamblador representativo,
    \item nombre de los ficheros RAM usados en la entrega,
    \item comentario del test,
    \item cronograma y/o captura del waveform si ayuda a justificar el comportamiento.
\end{itemize}

% -------------------------------------------------
\subsection{RF1 -- JAL y RET}
% -------------------------------------------------

\subsubsection*{Descripción}
Incluir las instrucciones \codigo{jal} y \codigo{ret} ya trabajadas en prácticas en el MIPS del proyecto.

\subsubsection*{¿Cómo?}
Explicar brevemente:
\begin{itemize}[leftmargin=1.2cm]
    \item los cambios realizados en el MIPS,
    \item cómo se calcula la dirección de salto,
    \item cómo se guarda la dirección de retorno,
    \item qué señales de control se activan,
    \item y cómo se actualiza el PC.
\end{itemize}

\vspace{0.3cm}
\noindent\pendiente

\subsubsection*{Verificación}
Comprobar su correcto funcionamiento en un banco de pruebas.

\paragraph{Título del test:} \pendiente

\paragraph{Ficheros RAM usados:}
\begin{itemize}[leftmargin=1.2cm]
    \item RAM de instrucciones: \pendiente
    \item RAM de datos: \pendiente
\end{itemize}

\paragraph{Código ensamblador representativo}
\begin{lstlisting}
# Pegar aquí el código del test de jal / ret
\end{lstlisting}

\paragraph{Comentario del test}
Explicar qué se comprueba y por qué demuestra el funcionamiento correcto.

\vspace{0.3cm}
\noindent\pendiente

% -------------------------------------------------
\subsection{RF2 -- Anticipación de operandos}
% -------------------------------------------------

\subsubsection*{Descripción}
El procesador es capaz de anticipar los operandos para las instrucciones \codigo{lw}, \codigo{sw} y aritméticas a distancia 1
(siempre que el dato a anticipar no venga de un \codigo{lw}) y a distancia 2.

\subsubsection*{¿Cómo?}
Explicar brevemente la gestión de anticipación de operandos:
\begin{itemize}[leftmargin=1.2cm]
    \item dónde está la unidad de forwarding,
    \item qué señales observa,
    \item qué multiplexores controla,
    \item qué casos se contemplan para \codigo{rs} y \codigo{rt}.
\end{itemize}

\vspace{0.3cm}
\noindent\pendiente

\subsubsection*{Verificación}
En la plantilla nueva se citan ejemplos del banco de pruebas \codigo{test\_IRQ}, como:
\begin{itemize}[leftmargin=1.2cm]
    \item de \codigo{LW R1, 8(R0)} a \codigo{ADD R31, R1, R31}: anticipación de \codigo{rs} a distancia 2,
    \item de \codigo{SUB R31, R31, R1} a \codigo{LW R1, 0(R31)}: anticipación de \codigo{rs} a distancia 1,
    \item de \codigo{ADD R2, R1, R2} a \codigo{SW R2, 7008(R6)}.
\end{itemize}

\paragraph{Código ensamblador representativo}
\begin{lstlisting}
# Pegar aquí casos representativos de forwarding
\end{lstlisting}

\paragraph{Comentario del test}
Completar con otros casos representativos y explicar por qué el valor correcto llega sin parada cuando procede.

\vspace{0.3cm}
\noindent\pendiente

% -------------------------------------------------
\subsection{RF3 -- Riesgos de datos}
% -------------------------------------------------

\subsubsection*{Descripción}
El procesador es capaz de detener la ejecución para evitar los riesgos causados por dependencias en instrucciones \codigo{lw}-uso
(distancia 1), \codigo{beq} o WRO (distancias 1 o 2).

\subsubsection*{¿Cómo?}
Explicar brevemente la gestión de detenciones:
\begin{itemize}[leftmargin=1.2cm]
    \item señales usadas por la unidad de detección,
    \item instrucciones que obligan a detener,
    \item detención de IF/ID,
    \item invalidación hacia EX.
\end{itemize}

\vspace{0.3cm}
\noindent\pendiente

\subsubsection*{Verificación}
En la plantilla nueva se citan estos casos en \codigo{test\_IRQ}:
\begin{itemize}[leftmargin=1.2cm]
    \item \textbf{Test 1:} 1 ciclo de detención por dependencia a distancia 2 entre \codigo{ADD R1, R1, R1} y \codigo{beq R1, R1, main},
    \item \textbf{Test 2:} 1 ciclo de detención por dependencia \codigo{lw}-uso, por ejemplo entre \codigo{LW R1, 8(R0)} y \codigo{SW R1, 7004(R6)},
    \item \textbf{Test 3:} no se para entre \codigo{SW R1, 7004(R6)} y \codigo{ADD R1, R1, R1}.
\end{itemize}

\paragraph{Código ensamblador representativo}
\begin{lstlisting}
# Pegar aquí pruebas de riesgos de datos
\end{lstlisting}

\paragraph{Comentario del test}
Completar con otros casos representativos e indicar cuántos ciclos de parada aparecen y por qué.

\vspace{0.3cm}
\noindent\pendiente

% -------------------------------------------------
\subsection{RF4 -- Riesgos de control}
% -------------------------------------------------

\subsubsection*{Descripción}
El procesador es capaz de eliminar los riesgos de control causados por las instrucciones de salto: \codigo{beq} y \codigo{rte}.

\subsubsection*{¿Cómo?}
Explicar brevemente:
\begin{itemize}[leftmargin=1.2cm]
    \item cómo se gestiona el PC,
    \item cuándo se anulan instrucciones,
    \item qué señales de control intervienen.
\end{itemize}

\vspace{0.3cm}
\noindent\pendiente

\subsubsection*{Verificación}
En el banco de pruebas \codigo{test\_IRQ} se verifican varios saltos tomados en los que sólo se ejecutan las instrucciones correctas:
\begin{itemize}[leftmargin=1.2cm]
    \item \codigo{beq R1, R1, main},
    \item \codigo{rte},
    \item un ejemplo en el que no se salta,
    \item verificar también \codigo{jal} y \codigo{ret}.
\end{itemize}

\paragraph{Código ensamblador representativo}
\begin{lstlisting}
# Pegar aquí pruebas de riesgos de control
\end{lstlisting}

\paragraph{Comentario del test}
Indicar qué instrucciones se anulan y cuáles llegan a ejecutarse realmente.

\vspace{0.3cm}
\noindent\pendiente

% -------------------------------------------------
\subsection{RF5 -- Riesgos estructurales}
% -------------------------------------------------

\subsubsection*{Descripción}
El procesador es capaz de detener la ejecución para evitar riesgos estructurales causados por las instrucciones \codigo{mac} y \codigo{mac\_ini}.

\subsubsection*{¿Cómo?}
Explicar brevemente la gestión de riesgos estructurales:
\begin{itemize}[leftmargin=1.2cm]
    \item qué etapas se detienen,
    \item qué señal de control se emplea,
    \item cómo se coordina con el resto de detenciones.
\end{itemize}

\vspace{0.3cm}
\noindent\pendiente

\subsubsection*{Verificación}
En el banco de pruebas \codigo{testbench\_2} se ejecutan varias instrucciones \codigo{mac} y \codigo{mac\_ini}, incluyendo:
\begin{itemize}[leftmargin=1.2cm]
    \item \codigo{mac\_ini R10, R1, R5},
    \item \codigo{mac R10, R2, R6}.
\end{itemize}

\paragraph{Código ensamblador representativo}
\begin{lstlisting}
# Pegar aquí test de riesgos estructurales con MAC / MAC_ini
\end{lstlisting}

\paragraph{Comentario del test}
Explicar la detención estructural y el resultado final.

\vspace{0.3cm}
\noindent\pendiente

% -------------------------------------------------
\subsection{RF6 -- Cambio a modo excepción}
% -------------------------------------------------

\subsubsection*{Descripción}
Si el procesador está en modo usuario y recibe \codigo{IRQ}, \codigo{ABORT} o \codigo{UNDEF}, pasa al modo correspondiente y ejecuta la rutina que indica la tabla de vectores de excepción.
Si está en modo excepción, ignora nuevas señales hasta volver a modo usuario.

\subsubsection*{¿Cómo?}
Explicar brevemente la gestión del PC al aceptar la excepción:
\begin{itemize}[leftmargin=1.2cm]
    \item detección de \codigo{IRQ}, \codigo{ABORT} y \codigo{UNDEF},
    \item selección de rutina,
    \item actualización de PC,
    \item cambio de modo.
\end{itemize}

\vspace{0.3cm}
\noindent\pendiente

\subsubsection*{Verificación}

\paragraph{IRQ}
Banco de pruebas \codigo{test\_IRQ} con diversas IRQs. Se atienden si y sólo si el procesador está en modo usuario.

\begin{lstlisting}
# Pegar aquí test de IRQ
\end{lstlisting}

\paragraph{Data Abort}
En la memoria de instrucciones se incluyen pruebas con:
\begin{itemize}[leftmargin=1.2cm]
    \item acceso no alineado,
    \item acceso a dirección fuera de rango.
\end{itemize}
En ambos casos la ejecución salta a la rutina Data Abort.

\begin{lstlisting}
# Pegar aquí test de Data Abort
\end{lstlisting}

\paragraph{Undef}
Se ejecuta una instrucción con código de operación desconocido y la ejecución salta a la rutina UNDEF.

\begin{lstlisting}
# Pegar aquí test de instrucción indefinida
\end{lstlisting}

\paragraph{Comentario del test}
Explicar cómo se observa en simulación la entrada al modo excepción y la rutina ejecutada.

\vspace{0.3cm}
\noindent\pendiente

% -------------------------------------------------
\subsection{RF7 -- Retorno a modo usuario con RTE}
% -------------------------------------------------

\subsubsection*{Descripción}
Al terminar la excepción se vuelve a la ejecución original con la instrucción \codigo{rte}.

\subsubsection*{¿Cómo?}
Explicar brevemente:
\begin{itemize}[leftmargin=1.2cm]
    \item cómo se restaura el PC,
    \item cómo se vuelve a modo usuario,
    \item cómo se conserva el contexto,
    \item cómo se usa el esquema de prólogo/epílogo.
\end{itemize}

\vspace{0.3cm}
\noindent\pendiente

\subsubsection*{Verificación}
Sólo es necesario verificarlo para las IRQs. El código debe funcionar aunque se interrumpa repetidas veces.
En \codigo{test\_IRQ} se comprueba que se retorna a la ejecución sin alterarla siempre que se guarden y restauren registros en pila, usando \codigo{R31} como SP.

\begin{lstlisting}
# Pegar aquí test de retorno con rte
\end{lstlisting}

\paragraph{Comentario del test}
Describir cómo se comprueba que el programa continúa correctamente tras la interrupción.

\vspace{0.3cm}
\noindent\pendiente

% -------------------------------------------------
\subsection{RF8 -- ALU multiciclo}
% -------------------------------------------------

\subsubsection*{Descripción}
El procesador es capaz de ejecutar las instrucciones \codigo{mac} y \codigo{mac\_ini} en 3 ciclos y gestionar el estado del registro \codigo{ACC} ante el tratamiento de excepciones.

\subsubsection*{¿Cómo?}
Explicar brevemente la máquina de estados diseñada:
\begin{itemize}[leftmargin=1.2cm]
    \item estados,
    \item entradas y salidas de cada estado,
    \item número de ciclos,
    \item gestión hardware o software del registro \codigo{ACC} al tratar excepciones.
\end{itemize}

\vspace{0.3cm}
\noindent\pendiente

\subsubsection*{Verificación}
El banco de pruebas \codigo{testbench\_2} incluye varias instrucciones \codigo{mac} y \codigo{mac\_ini}.

Además, conviene diseñar un banco de pruebas donde tanto la función principal como la RTI empleen instrucciones MAC y verificar el estado del registro \codigo{ACC}.

\begin{lstlisting}
# Pegar aquí test de ALU multiciclo / MAC
\end{lstlisting}

\paragraph{Comentario del test}
Explicar el funcionamiento en 3 ciclos y el tratamiento del registro \codigo{ACC}.

\vspace{0.3cm}
\noindent\pendiente

% -------------------------------------------------
\subsection{RF9 -- Contadores}
% -------------------------------------------------

\subsubsection*{Descripción}
Los contadores dan información de la ejecución del código.

\subsubsection*{¿Cómo?}
Explicar brevemente la gestión de los contadores:
\begin{itemize}[leftmargin=1.2cm]
    \item cuándo incrementa cada contador,
    \item qué señales los habilitan,
    \item cómo se evita contar una parada varias veces,
    \item jerarquía entre causas de parada.
\end{itemize}

\vspace{0.3cm}
\noindent\pendiente

\subsubsection*{Verificación}
En el banco de pruebas \codigo{test\_IRQ} se ha comprobado que:
\begin{itemize}[leftmargin=1.2cm]
    \item \codigo{Ins} incrementa por cada instrucción válida que hace WB,
    \item \codigo{Data\_stalls} contabiliza los ciclos debidos a riesgos de datos,
    \item \codigo{Control\_stalls} incrementa un ciclo por salto tomado,
    \item \codigo{Exception\_accepted} contabiliza las excepciones aceptadas,
    \item \codigo{Memory\_stalls} incrementa cuando memoria no está \codigo{ready}.
\end{itemize}

\paragraph{Jerarquía de paradas}
Importante: las paradas deben incrementar sólo un contador.
Si hay varios motivos, se cuenta el de mayor jerarquía:
\[
\texttt{memory stall} > \texttt{data stall} > \texttt{control stall}
\]

\begin{lstlisting}
# Pegar aquí programa o test para verificar contadores
\end{lstlisting}

\paragraph{Comentario del test}
Explicar la relación entre el comportamiento observado y los valores leídos en los contadores.

\vspace{0.3cm}
\noindent\pendiente

% -------------------------------------------------
\section{Cuantificación de horas dedicadas}
% -------------------------------------------------

Indicar horas por tarea (y persona en el caso de parejas).

\begin{center}
\begin{tabular}{p{9cm}c}
\toprule
\textbf{Tarea} & \textbf{Horas} \\
\midrule
Estudio del MIPS, VHDL, entorno, instalación\ldots & \pendiente \\
Adición de las nuevas instrucciones & \pendiente \\
Gestión de excepciones & \pendiente \\
Gestión de riesgos & \pendiente \\
Depuración, verificación y programas de prueba & \pendiente \\
Memoria & \pendiente \\
\midrule
\textbf{Total} & \pendiente \\
\bottomrule
\end{tabular}
\end{center}

% -------------------------------------------------
\section{Conclusiones y autoevaluación}
% -------------------------------------------------

Responder aquí a cuestiones como:
\begin{itemize}[leftmargin=1.2cm]
    \item ¿Para qué ha servido vuestro trabajo?
    \item ¿Qué os ha parecido?
    \item ¿Qué calificación os otorgáis a vosotros mismos?
\end{itemize}

\vspace{0.5cm}
\noindent\pendiente

% -------------------------------------------------
\section{Agradecimientos}
% -------------------------------------------------

¿Os ha ayudado algún compañer@? Podéis indicarlo aquí.

\vspace{0.5cm}
\noindent\pendiente

\end{document}